\documentclass[11pt,a4paper]{article}
\usepackage[margin=1.8cm,top=2cm,bottom=2.2cm]{geometry}
\usepackage{tabularx}
\usepackage{booktabs}
\usepackage{qrcode}
\usepackage{fancyhdr}
\usepackage{lastpage}
\usepackage{xcolor}

\definecolor{primary}{HTML}{1E293B}
\definecolor{muted}{HTML}{64748B}

\pagestyle{fancy}
\fancyhf{}
\renewcommand{\headrulewidth}{0pt}
\fancyfoot[L]{\footnotesize\color{muted} Sistem Administrasi Kawasan Pergudangan --- Dicetak otomatis pada {{PRINTED_AT}}}
\fancyfoot[R]{\footnotesize\color{muted} Halaman \thepage\ dari \pageref{LastPage}}

\begin{document}

\noindent
\begin{tabularx}{\textwidth}{@{} X r @{}}
  {\Large\bfseries\color{primary} KAWASAN PERGUDANGAN TERPADU} \newline
  {\normalsize\bfseries SURAT IZIN MASUK KAWASAN (GATE PASS)} \newline
  {\small\color{muted} Nomor Dokumen: \textbf{{{DOCUMENT_NUMBER}}} \quad|\quad Tiket Antrean: \textbf{{{TICKET_NUMBER}}}} &
  \qrcode[height=2.2cm]{{{QR_PAYLOAD}}}
\end{tabularx}

\vspace{0.3cm}
\noindent\rule{\textwidth}{1.2pt}
\vspace{0.3cm}

\section*{1. Informasi Kunjungan \& Armada}
\begin{tabularx}{\textwidth}{@{} p{4.6cm} X @{}}
  \toprule
  \textbf{Waktu Gate-In} & {{GATE_IN_AT}} \\
  \textbf{Gudang Tujuan} & {{WAREHOUSE_CODE}} --- {{WAREHOUSE_NAME}} \\
  \textbf{Arah Pergerakan} & \textbf{{{DIRECTION}}} \\
  \textbf{Kategori Angkutan} & \textbf{{{CATEGORY}}} \\
  \textbf{Nomor Polisi Armada} & \textbf{{{PLATE_NUMBER}}} ({{VEHICLE_TYPE}}) \\
  \textbf{Perusahaan Ekspedisi} & {{COMPANY_NAME}} \\
  \textbf{Kapasitas Maksimal} & {{MAX_WEIGHT_KG}} kg \\
  \textbf{Nama Pengemudi} & {{DRIVER_NAME}} \\
  \textbf{Nomor SIM Pengemudi} & {{DRIVER_LICENSE_NO}} \\
  \bottomrule
\end{tabularx}

\vspace{0.4cm}
\section*{2. Rincian Kontainer (Khusus Trailer / Tronton)}
{{CONTAINER_SECTION}}

\vspace{0.5cm}
\noindent\textbf{Catatan Operasional:}
\begin{itemize}
  \item Pengemudi wajib mematuhi batas kecepatan maksimal 20 km/jam di dalam kawasan pergudangan.
  \item Armada hanya diperbolehkan merapat pada \textit{loading dock} yang telah dialokasikan oleh Admin Gudang.
  \item Dokumen ini wajib diserahkan kembali di Pos Gerbang Keluar bersama BASTB atau Surat Jalan resmi.
\end{itemize}

\vspace{1.2cm}
\noindent
\begin{tabularx}{\textwidth}{@{} X X @{}}
  \centering
  Petugas Pos Gerbang (Security), \vspace{1.8cm} \newline
  \textbf{(\underline{\hspace{4.5cm}})} \newline
  {\small {{ISSUED_BY_NAME}}} &
  \centering
  Pengemudi Armada, \vspace{1.8cm} \newline
  \textbf{(\underline{\hspace{4.5cm}})} \newline
  {\small {{DRIVER_NAME}}}
\end{tabularx}

\end{document}
